\documentclass[12pt]{article}
\usepackage[utf8]{inputenc}
\usepackage[T1]{fontenc}
\usepackage[margin=1in]{geometry}
\usepackage{setspace}
\usepackage{xcolor}
\usepackage[colorlinks=true, urlcolor=blue!60!black, linkcolor=blue!60!black]{hyperref}
\setlength{\parindent}{2em}
\setlength{\parskip}{1.25em}
\renewcommand{\baselinestretch}{1.0}
\begin{document}
\begin{center}
{\LARGE \textit{Moody v.\ NetChoice, LLC}, 603 U.S.\ 707 (2024)\par}
\vspace{0.8em}
{\large Case Brief --- Four-Page Edition\par}
\end{center}
\vspace{1em}

\section{Memory Jogger}

The Court unanimously vacated both lower-court judgments because neither applied the demanding test for facial First Amendment challenges; a majority (Kagan, J.) added that curating third-party posts in a feed is protected editorial expression that a State may not override to rebalance the marketplace of ideas (\href{https://supreme.justia.com/cases/federal/us/603/22-277/}{Justia}).

\section{Detailed Case Facts}

In 2021, amid concerns that large technology companies were suppressing conservative viewpoints, Florida and Texas restricted social-media content moderation. \textbf{Florida S.B.~7072} covers platforms with over \$100 million in revenue or 100 million monthly users; it bars ``censoring,'' deprioritizing, or labeling posts by or about candidates or ``journalistic enterprises,'' requires ``consistent'' moderation, limits rule changes, mandates a chronological opt-out, and requires a ``thorough rationale'' for removals. \textbf{Texas H.B.~20} covers platforms with over 50 million U.S.\ users and forbids ``censoring'' (blocking, removing, demonetizing, or de-boosting) users or expression based on viewpoint, with explanations and a 14-day appeal.

\section{Procedural History}

NetChoice and the Computer \& Communications Industry Association, representing Facebook, YouTube, X, and others, brought pre-enforcement \emph{facial} challenges. District courts enjoined both laws. The Eleventh Circuit largely affirmed, treating moderation as protected editorial discretion and the explanation rules as unduly burdensome under \textit{Zauderer}; the Fifth Circuit reversed, calling moderation ``not speech'' and upholding Texas's interest in ``protecting a diversity of ideas.''

\section{Judicial Votes}

Vacated and remanded, \textbf{9--0} (decided July 1, 2024).
\begin{itemize}
\item \textbf{Majority:} Justice Kagan, joined by Chief Justice Roberts and Justices Sotomayor, Kavanaugh, and Barrett, and by Justice Jackson as to Parts I, II, and III--A.
\item \textbf{Concurrences:} Justice Barrett; Justice Jackson (in part and in the judgment); Justice Thomas (in the judgment); Justice Alito (in the judgment), joined by Thomas and Gorsuch.
\item \textbf{Dissents:} None.
\end{itemize}

\section{Holding}

Neither circuit conducted the required facial inquiry: define each statute's full scope and ask whether a ``substantial number of [the law's] applications are unconstitutional, judged in relation to the statute's plainly legitimate sweep.'' To guide remand, a majority held that curating third-party posts on feeds like Facebook's News Feed and YouTube's homepage is protected expressive activity under \textit{Tornillo}, \textit{PG\&E}, \textit{Turner}, and \textit{Hurley}, and that an interest in correcting perceived ideological bias fails even intermediate scrutiny.

\section{Analysis of Opinions}

\textbf{Majority (Kagan): Facial challenges.} Facial challenges are disfavored. Even under the First Amendment's relaxed standard, a law must prohibit ``a substantial amount of protected speech relative to its plainly legitimate sweep.'' The courts below treated the suits as challenges to two feeds, without asking whether the laws reach direct messaging, Gmail, Etsy, Venmo, or Uber, or comparing applications. The Court is ``a court of review, not of first view.''

\textbf{Majority: Editorial discretion.} Compiling, filtering, prioritizing, and labeling third-party speech is protected expression; a host that excludes only a few messages keeps protection; and government may not override editorial choices to ``better balance'' the marketplace. \textit{PruneYard} and \textit{FAIR} involved hosts not engaged in their own expression. Feeds applying Community Standards to hate speech or misinformation produce a distinctive product, which Texas's mandate alters. Texas's aim ``to correct the mix of speech'' is related to suppressing expression and fails \textit{O'Brien}; under \textit{Buckley}, the State may not ``restrict the speech of some elements of our society in order to enhance the relative voice of others.''

\textbf{Barrett.} Facial suits ask courts ``to bite off more than [they] can chew''; AI moderation and foreign ownership remain open. \textbf{Jackson.} She declined to preview the merits. \textbf{Thomas.} He would revisit \textit{Zauderer}, stressed common-carrier doctrine, and argued facial challenges exceed Article~III. \textbf{Alito.} Part III is ``nonbinding dicta''; \textit{PruneYard}, \textit{FAIR}, and common-carrier history deserve serious attention, and AI-driven platforms differ from newspaper editors.

\textbf{Later.} On remand the Fifth Circuit ordered ``thorough discovery,'' and lower courts have used \textit{Moody} to vacate injunctions against youth social-media laws.

\section{Examples: Future Applications}

\begin{enumerate}
\item \textbf{Same-Side: A Climate-Viewpoint Mandate.} A facial suit addressing only one curated feed fails, but as applied to that feed, the rebalancing law is invalid.
\item \textbf{Same-Side: Must-Carry for Candidates.} Fining a video platform for labeling a candidate's false voting-date claim overrides protected editorial judgment, as in \textit{Tornillo}.
\item \textbf{Opposite-Side: Payment Apps and Messaging.} Barring payment apps from blocking transfers by political affiliation regulates conduits that compile no expressive product.
\item \textbf{Opposite-Side: Ride-Share Drivers.} Forbidding deactivation of drivers for off-app political speech is economic regulation, as in \textit{PruneYard} and \textit{FAIR}.
\item \textbf{Fence-Sitter: Engagement-Only Algorithms.} A feed ranked solely to maximize engagement, with no content standards, is the case the Court reserved.
\end{enumerate}

\section{Critique}

\textbf{Holding versus dicta.} Five Justices joined Part III's defense of editorial discretion, but three called it ``nonbinding dicta'' and Jackson abstained. \textbf{The algorithm gap.} Footnote~6 and Barrett leave open whether engagement algorithms or generative-AI moderators exercise editorial judgment. \textbf{Procedural asymmetry.} Mapping every application, from Gmail to Uber, raises the cost of facial relief and pushes platforms toward as-applied suits.

\section{Key Quotations}

\begin{quote}
``Deciding on the third-party speech that will be included in or excluded from a compilation---and then organizing and presenting the included items---is expressive activity of its own. And that activity results in a distinctive expressive product.''
\end{quote}
\noindent --- Justice Kagan, for the Court, 603 U.S.\ at 731.

\begin{quote}
``On the spectrum of dangers to free expression, there are few greater than allowing the government to change the speech of private actors in order to achieve its own conception of speech nirvana.''
\end{quote}
\noindent --- Justice Kagan, for the Court, 603 U.S.\ at 741--42.

\begin{quote}
``In sum, there is much work to do below on both these cases, given the facial nature of NetChoice's challenges. But that work must be done consistent with the First Amendment, which does not go on leave when social media are involved.''
\end{quote}
\noindent --- Justice Kagan, for the Court, 603 U.S.\ at 719.

\end{document}
